\section{The axial coordinate and the cascade orientation}
\label{sec:coordinate}

\subsection{$z$ becomes a depth}

With the return no longer exchanging, the developed length of a heat-transfer
path is $L_t = \Lw$ rather than $2\Lw$, and the map onto depth is one-to-one.
This single change is what admits a depth-dependent formation.

On charge the water reaches the bottom of the hole through the downcomer
before it has given up any heat, so the finned tubes are entered at the
bottom: the solver coordinate $z = 0$ is the bottom of the well and depth
decreases with $z$,
\begin{equation}
\text{depth}(z) = \Lw - z .
\end{equation}
On discharge the flow reverses --- cold water descends the insulated
downcomer, enters the tubes at the bottom and rises --- so the march is run
with the segment index reversed, and the formation profile is reversed with
it, by exactly the same rule already applied to the melting-temperature
profile. A mirroring error here would place the shielded shallow resistance
against the deep rock, so the orientation is asserted in the parameter
validation rather than left to a comment.

\subsection{The cascade necessarily runs hot at the bottom}

Because the charge water arrives at the bottom undiminished, the hottest
melting layer is at the bottom and the coldest at the top. This is not a free
choice: the geometry determines it.

It is also the favourable orientation. The formation is coldest at the surface
and hottest at depth, so a hot-at-top cascade --- the THUMS arrangement ---
opposes the gradient at both ends. The net driving difference is identical for
either orientation, since the depth-mean of $\Tm(z)$ is unchanged, but the
distribution is not:

\begin{table}[h]
\centering
\begin{tabular}{lrrr}
\toprule
at \SI{8000}{ft}, $\nabla T = \SI{56.5}{\kelvin\per\kilo\metre}$
 & $\Delta T$ at top & $\Delta T$ at bottom & cannot freeze \\
\midrule
hot-at-top (THUMS)     & $+132$\,K & $-55$\,K & \SI{29}{\percent} \\
hot-at-bottom (BOREAS) & $+83$\,K  & $-6$\,K  & \SI{7}{\percent}  \\
\bottomrule
\end{tabular}
\caption{Material held above its melting temperature by the surrounding rock
cannot solidify on discharge, and that capacity is dead --- drilled and paid
for, and unable to cycle. Both orientations have a depth-mean driving
difference of \SI{38.7}{\kelvin}; only the distribution differs.}
\end{table}

